% !TeX program = xelatex
% ============================================================
% rp/rp.tex — Руководство программиста
% ГОСТ 19.504-79 «Руководство программиста. Требования к
% содержанию и оформлению»
% ============================================================
\documentclass[12pt,a4paper]{extarticle}

% 1. Подключение общей преамбулы (пакеты)
\input{../common/preamble}

% 2. Определение переменных для конкретного документа
\newcommand{\docname}{Руководство программиста}
\newcommand{\doccode}{\hseDocCodeBase\ 33 01-1}
\newcommand{\doccodelu}{\hseDocCodeBase\ 33 01-1-ЛУ}

% 3. Подключение общих команд (проект, студент, руководитель)
\input{../common/commands}

% 4. Подключение стилей (колонтитулы, заголовки)
\input{../common/styles}

\begin{document}
\sloppy

% 5. Генерация титульных листов по шаблону
\input{../common/title_template}

% ============================================================
% АННОТАЦИЯ
% ============================================================
\section*{АННОТАЦИЯ}
\addcontentsline{toc}{section}{АННОТАЦИЯ}

Настоящий документ «\hseProjectName. Руководство программиста» содержит сведения, необходимые для разработки, настройки, запуска, проверки и сопровождения программного комплекса. В документе описаны назначение и условия применения программы, характеристики архитектуры, порядок обращения к программе, организация входных и выходных данных, а также диагностические сообщения, формируемые системой.

% Подсказка: уточните, для кого предназначено руководство (какие сервисы,
% библиотеки и конфигурации разрабатывают его читатели).
Руководство предназначено для программистов и инженеров, выполняющих разработку сервисов, общих библиотек, программных контрактов и инфраструктурных конфигураций проекта.

Документ разработан в соответствии с требованиями ГОСТ 19.504-79\cite{gost19504}.

\clearpage

% ============================================================
% СПИСОК ИСПОЛЬЗУЕМЫХ ГОСТ
% ============================================================
\noindent
Настоящий документ разработан в~соответствии с~требованиями:

\begin{enumerate}[label=\arabic*), leftmargin=1.25cm]
  \item ГОСТ 19.101-77 Виды программ и~программных документов\cite{gost19101}.
  \item ГОСТ 19.103-77 Обозначения программ и~программных документов\cite{gost19103}.
  \item ГОСТ 19.104-78 Основные надписи\cite{gost19104}.
  \item ГОСТ 19.105-78 Общие требования к~программным документам\cite{gost19105}.
  \item ГОСТ 19.106-78 Требования к~программным документам, выполненным
        печатным способом\cite{gost19106}.
  \item ГОСТ 19.504-79 Руководство программиста. Требования
        к~содержанию и~оформлению\cite{gost19504}.
\end{enumerate}

Изменения к~данному документу оформляются согласно ГОСТ\,19.604-78\cite{gost19604}.

\clearpage

% ============================================================
% СОДЕРЖАНИЕ
% ============================================================
\hypersetup{linkcolor=black}
\tableofcontents
\hypersetup{linkcolor=blue}
\thispagestyle{fancy}
\clearpage

% ============================================================
% 1. НАЗНАЧЕНИЕ И УСЛОВИЯ ПРИМЕНЕНИЯ ПРОГРАММЫ
% ============================================================
\section{НАЗНАЧЕНИЕ И УСЛОВИЯ ПРИМЕНЕНИЯ ПРОГРАММЫ}

\subsection{Назначение программы}

% Подсказка: опишите назначение программы и круг решаемых ею задач.
\begin{hseExample}
Программа «\projectname» предназначена для \hseFill{укажите назначение программы}. Система предоставляет \hseFill{категории пользователей} программные интерфейсы для \hseFill{перечень решаемых задач}.

Программа применяется как \hseFill{роль программы в общей системе} и как основа для последующего развития \hseFill{направления развития функциональности}.
\end{hseExample}

\subsection{Функции, выполняемые программой}

% Подсказка: перечислите основные функции программы (список через itemize).
\begin{hseExample}
Программный комплекс выполняет следующие основные функции:

\begin{itemize}[leftmargin=1.25cm]
    \item \hseFill{первая основная функция программы};
    \item \hseFill{вторая основная функция программы};
    \item \hseFill{третья основная функция программы};
    \item сбор метрик, трассировок и структурированных журналов.
\end{itemize}
\end{hseExample}

\subsection{Условия, необходимые для выполнения программы}

% Подсказка: укажите окружение, в котором запускается программа
% (контейнеры, ОС, среда разработки и т. п.).
\begin{hseExample}
Программа рассчитана на запуск в \hseFill{окружение выполнения программы}. Для локальной разработки используются \hseFill{инструменты локальной разработки}.
\end{hseExample}

\subsubsection{Требования к аппаратным средствам}

% Подсказка: минимальные/рекомендуемые требования к процессору, ОЗУ,
% дисковому пространству и сети.
\begin{hseExample}
Для локального запуска программы рекомендуется рабочая станция со следующими характеристиками:

\begin{itemize}[leftmargin=1.25cm]
    \item процессор: не менее \hseFill{число ядер и архитектура};
    \item оперативная память: не менее \hseFill{объем ОЗУ};
    \item дисковое пространство: не менее \hseFill{объем диска} свободного пространства для исходного кода, зависимостей и данных;
    \item сетевой доступ к \hseFill{репозиторий и внешние ресурсы}.
\end{itemize}

Для промышленного запуска компоненты должны размещаться в \hseFill{целевая среда развертывания} с выделенными ресурсами для \hseFill{СУБД и инфраструктурные сервисы}.
\end{hseExample}

\subsubsection{Требования к программным средствам}

% Подсказка: перечислите ОС, рантаймы, СУБД, инструменты сборки и прочее ПО.
\begin{hseExample}
Для разработки и запуска программы требуются следующие программные средства:

\begin{itemize}[leftmargin=1.25cm]
    \item операционная система \hseFill{поддерживаемые ОС};
    \item Git для получения исходного кода;
    \item \hseFill{средства контейнеризации или виртуализации};
    \item \hseFill{язык программирования и его версия};
    \item \hseFill{СУБД и инфраструктурные зависимости};
    \item \hseFill{дополнительные инструменты разработки}.
\end{itemize}
\end{hseExample}

\subsubsection{Требования к персоналу}

% Подсказка: укажите квалификацию программиста (языки, технологии, навыки).
\begin{hseExample}
Программист, выполняющий разработку и сопровождение программы, должен обладать практическими навыками работы с \hseFill{языки, фреймворки и инструменты} и системой контроля версий Git. Для изменения \hseFill{смежный слой системы} дополнительно требуются навыки работы с \hseFill{дополнительные технологии и инструменты}.
\end{hseExample}

\clearpage

% ============================================================
% 2. ХАРАКТЕРИСТИКИ ПРОГРАММЫ
% ============================================================
\section{ХАРАКТЕРИСТИКИ ПРОГРАММЫ}

\subsection{Описание основных характеристик программы}

% Подсказка: опишите архитектуру, используемые подходы и технологии.
\begin{hseExample}
Программа «\projectname» реализована как \hseFill{тип программного комплекса}. Серверная часть построена на \hseFill{языки и фреймворки} и использует \hseFill{архитектурный стиль}. Взаимодействие между компонентами выполняется через \hseFill{протоколы взаимодействия}.

В системе используются следующие архитектурные подходы:

\begin{itemize}[leftmargin=1.25cm]
    \item разделение компонентов по зонам ответственности;
    \item \hseFill{второй архитектурный подход};
    \item \hseFill{третий архитектурный подход};
    \item миграции схем данных через \hseFill{инструмент миграций}.
\end{itemize}
\end{hseExample}

\subsection{Режим работы программы}

% Подсказка: укажите режим работы (непрерывный/пакетный/интерактивный) и
% допустимые сценарии запуска.
\begin{hseExample}
В штатном режиме программа работает непрерывно. \hseFill{опишите работу компонентов в штатном режиме}.

Для разработки допускаются следующие режимы запуска:

\begin{itemize}[leftmargin=1.25cm]
    \item запуск только инфраструктурных зависимостей;
    \item запуск выбранного компонента программы;
    \item запуск всех компонентов одновременно;
    \item запуск тестового окружения для интеграционных проверок.
\end{itemize}
\end{hseExample}

\subsection{Состав и назначение компонентов}

% Подсказка: перечислите основные компоненты программы и их назначение.
\begin{hseExample}
Состав и назначение основных компонентов программы приведены в таблице~\ref{tab:rp_components}.

% Таблица набрана через \captionof (не \begin{table}[H]): плавающие
% окружения внутри жёлтого блока-образца недопустимы.
\begin{center}
\small
\captionof{table}{Основные компоненты программы}
\label{tab:rp_components}
\begin{tabularx}{\textwidth}{|>{\raggedright\arraybackslash}p{0.34\textwidth}|X|}
\hline
\textbf{Компонент} & \textbf{Назначение} \\
\hline
\hseFill{имя компонента} & \hseFill{назначение компонента}. \\
\hline
\hseFill{имя компонента} & \hseFill{назначение компонента}. \\
\hline
\end{tabularx}
\end{center}
\end{hseExample}

\subsection{Средства контроля правильности выполнения}

% Подсказка: опишите средства контроля корректности работы (тесты,
% линтеры, health-checks, логи, метрики, трассировки).
\begin{hseExample}
Контроль корректности работы программы выполняется следующими средствами:

\begin{itemize}[leftmargin=1.25cm]
    \item автоматические unit- и integration-тесты через \hseFill{тестовый фреймворк};
    \item проверки форматирования, статического анализа и типов через \hseFill{линтеры и анализаторы};
    \item проверки работоспособности (health checks) компонентов программы;
    \item метрики, трассировки и структурированные журналы;
    \item \hseFill{дополнительные средства контроля}.
\end{itemize}
\end{hseExample}

\clearpage

% ============================================================
% 3. ОБРАЩЕНИЕ К ПРОГРАММЕ
% ============================================================
\section{ОБРАЩЕНИЕ К ПРОГРАММЕ}

\subsection{Получение исходного кода}

% Подсказка: укажите расположение репозитория и команды клонирования.
\begin{hseExample}
Исходный код программы расположен в удаленном Git-репозитории \hseFill{укажите git-хостинг и имя репозитория}.

Для получения исходного кода необходимо выполнить:

\begin{verbatim}
git clone <URL репозитория>
cd <папка проекта>
\end{verbatim}
\end{hseExample}

\subsection{Подготовка окружения разработки}

% Подсказка: опишите подготовку конфигураций и установку зависимостей.
\begin{hseExample}
Перед запуском программы необходимо установить системные зависимости, указанные в подразделе 1.3.2, и подготовить файлы переменных окружения. Примеры конфигураций \texttt{.env.example} расположены в \hseFill{каталоги с примерами конфигураций}.

Минимальный порядок подготовки:

\begin{enumerate}[label=\arabic*), leftmargin=1.25cm]
    \item скопировать примеры конфигураций в рабочие \texttt{.env}-файлы;
    \item установить зависимости командой \hseFill{команда установки зависимостей};
    \item \hseFill{дополнительные шаги подготовки окружения};
    \item применить миграции базы данных.
\end{enumerate}
\end{hseExample}

\subsection{Запуск программы}

% Подсказка: приведите команды запуска программы и её компонентов
% (инфраструктурные зависимости, отдельный компонент, всё окружение).
\begin{hseExample}
Для запуска инфраструктурных зависимостей используется команда:

\begin{verbatim}
<команда запуска инфраструктуры>
\end{verbatim}

Для запуска программы и всех её компонентов:

\begin{verbatim}
<команда запуска программы>
\end{verbatim}

Остановка локального окружения выполняется командой:

\begin{verbatim}
<команда остановки>
\end{verbatim}
\end{hseExample}

\subsection{Конфигурирование}

% Подсказка: опишите параметры конфигурации (переменные окружения,
% конфигурационные файлы) и их значения по умолчанию.
\begin{hseExample}
Параметры программы задаются переменными окружения и конфигурационными файлами. К основным параметрам относятся строки подключения к \hseFill{инфраструктурные зависимости}, параметры наблюдаемости, ключи сервисов и флаги запуска компонентов.
\end{hseExample}

\subsection{Выполнение проверок}

% Подсказка: приведите команды запуска тестов и других проверок.
\begin{hseExample}
Для запуска тестов, статических проверок и форматирования используются команды:

\begin{verbatim}
<команда запуска тестов>
<команда статических проверок>
<команда форматирования>
\end{verbatim}

Для smoke-проверки локально запущенного окружения:

\begin{verbatim}
<команда smoke-проверки>
\end{verbatim}
\end{hseExample}

\clearpage

% ============================================================
% 4. ВХОДНЫЕ И ВЫХОДНЫЕ ДАННЫЕ
% ============================================================
\section{ВХОДНЫЕ И ВЫХОДНЫЕ ДАННЫЕ}

\subsection{Организация входных данных}

% Подсказка: опишите источники, форматы и состав входных данных.
\begin{hseExample}
Входные данные программы поступают через \hseFill{интерфейсы и каналы входных данных}, переменные окружения, файлы конфигурации и миграционные данные. Основные форматы входных данных приведены в таблице~\ref{tab:rp_input_data}.

% Таблица набрана через \captionof (не \begin{table}[H]): плавающие
% окружения внутри жёлтого блока-образца недопустимы.
\begin{center}
\small
\captionof{table}{Входные данные программы}
\label{tab:rp_input_data}
\begin{tabularx}{\textwidth}{|>{\raggedright\arraybackslash}p{0.28\textwidth}|X|}
\hline
\textbf{Источник} & \textbf{Состав данных} \\
\hline
\hseFill{интерфейс API} & \hseFill{состав данных запросов}. \\
\hline
Переменные окружения & Строки подключения к \hseFill{инфраструктурные зависимости}, параметры наблюдаемости, ключи и флаги запуска компонентов. \\
\hline
Миграции БД & Скрипты \hseFill{инструмент миграций} для изменения структуры базы данных. \\
\hline
\end{tabularx}
\end{center}
\end{hseExample}

\subsection{Организация выходных данных}

% Подсказка: опишите каналы, форматы и состав выходных данных.
\begin{hseExample}
Выходные данные формируются в виде \hseFill{каналы выходных данных}, записей баз данных, метрик, трассировок и журналов. Основные виды выходной информации приведены в таблице~\ref{tab:rp_output_data}.

% Таблица набрана через \captionof (не \begin{table}[H]): плавающие
% окружения внутри жёлтого блока-образца недопустимы.
\begin{center}
\small
\captionof{table}{Выходные данные программы}
\label{tab:rp_output_data}
\begin{tabularx}{\textwidth}{|>{\raggedright\arraybackslash}p{0.28\textwidth}|X|}
\hline
\textbf{Канал} & \textbf{Состав данных} \\
\hline
\hseFill{канал API} & \hseFill{состав ответов и ошибок}. \\
\hline
База данных & Бизнес-данные программы, миграционная история и служебные структуры. \\
\hline
Observability & Метрики, трассировки, события ошибок и структурированные журналы. \\
\hline
\end{tabularx}
\end{center}
\end{hseExample}

\subsection{Ограничения и валидация данных}

% Подсказка: опишите уровни валидации и ограничения на данные.
\begin{hseExample}
Валидация входных данных выполняется на нескольких уровнях:

\begin{itemize}[leftmargin=1.25cm]
    \item \hseFill{схемы валидации} проверяют типы, обязательность и формат полей;
    \item доменный слой проверяет бизнес-инварианты: \hseFill{примеры бизнес-правил};
    \item ограничения и индексы базы данных защищают от нарушения целостности данных;
    \item \hseFill{дополнительные механизмы защиты данных}.
\end{itemize}
\end{hseExample}

\clearpage

% ============================================================
% 5. СООБЩЕНИЯ
% ============================================================
\section{СООБЩЕНИЯ}

% Подсказка: опишите типы формируемых сообщений и заполните таблицу
% диагностических сообщений (код/текст, причина, действия программиста).
\begin{hseExample}
Программа формирует сообщения нескольких типов: \hseFill{типы ответов и кодов ошибок}, записи журналов, сообщения проверок работоспособности и события мониторинга. Сообщения предназначены для программиста, оператора системы и вызывающих программных клиентов. Основные диагностические сообщения приведены в таблице~\ref{tab:rp_messages}.
\end{hseExample}

% longtable остаётся вне жёлтого блока-образца: разрывные таблицы внутри
% hseExample недопустимы.
\begin{small}
\begin{longtable}{|>{\raggedright\arraybackslash}p{0.19\textwidth}|>{\raggedright\arraybackslash}p{0.21\textwidth}|>{\raggedright\arraybackslash}p{0.48\textwidth}|}
\caption{Основные диагностические сообщения и действия программиста}
\label{tab:rp_messages}\\
\hline
\textbf{Сообщение/код} & \textbf{Причина} & \textbf{Действия программиста} \\
\hline
\endfirsthead
\hline
\textbf{Сообщение/код} & \textbf{Причина} & \textbf{Действия программиста} \\
\hline
\endhead
\hline
\endfoot
\texttt{400 Bad Request} & Некорректный формат входных данных или параметров запроса. & Проверить схему запроса, клиентскую валидацию и тесты граничных значений. \\
\hline
\texttt{404 Not Found} & Запрошенная сущность не найдена. & Проверить идентификатор, данные и корректность маршрута. \\
\hline
\texttt{503 Service Unavailable} & Недоступен зависимый сервис, база данных или внешний компонент. & Проверить health checks, переменные окружения, сетевые настройки и логи соответствующего сервиса. \\
\hline
\texttt{Health check not ready} & Компонент запущен, но не готов принимать трафик. & Проверить readiness endpoint, миграции и соединения с зависимостями. \\
\hline
\hseFill{сообщение или код} & \hseFill{причина} & \hseFill{действия программиста} \\
\hline
\end{longtable}
\end{small}

\clearpage
% Страницы терминов, источников и приложений печатаются без нижнего штампа
% (как в реальных комплектах) — только номер листа и код документа сверху.
\pagestyle{appendix}

% ============================================================
% ПЕРЕЧЕНЬ ТЕРМИНОВ И ОПРЕДЕЛЕНИЙ
% ============================================================
\section*{ПЕРЕЧЕНЬ ТЕРМИНОВ И ОПРЕДЕЛЕНИЙ}
\addcontentsline{toc}{section}{ПЕРЕЧЕНЬ ТЕРМИНОВ И ОПРЕДЕЛЕНИЙ}

% Подсказка: перечислите термины и сокращения, использованные в документе
% (в алфавитном порядке).
\begin{enumerate}[label=\arabic*., leftmargin=1.25cm]
  \item \textbf{\hseFill{термин}} --- \hseFill{определение термина}.
\end{enumerate}

\clearpage

% ============================================================
% СПИСОК ИСПОЛЬЗОВАННЫХ ИСТОЧНИКОВ
% ============================================================
\section*{СПИСОК ИСПОЛЬЗОВАННЫХ ИСТОЧНИКОВ}
\addcontentsline{toc}{section}{СПИСОК ИСПОЛЬЗОВАННЫХ ИСТОЧНИКОВ}

% Подсказка: добавьте ПЕРЕД ГОСТами источники по технологиям вашего
% проекта в алфавитном порядке, по образцу:
% \bibitem{fastapi}
% FastAPI [Электронный ресурс]. Режим доступа: \url{https://fastapi.tiangolo.com}, свободный. (дата обращения: 14.03.2026).
\renewcommand{\refname}{}
\begin{thebibliography}{99}
\vspace{-1cm}

\bibitem{gost19101}
ГОСТ 19.101-77: Виды программ и~программных документов. // Единая система программной документации.~--- М.: ИПК Издательство стандартов, 2001.

\bibitem{gost19103}
ГОСТ 19.103-77: Обозначения программ и~программных документов. // Единая система программной документации.~--- М.: ИПК Издательство стандартов, 2001.

\bibitem{gost19104}
ГОСТ 19.104-78: Основные надписи. // Единая система программной документации.~--- М.: ИПК Издательство стандартов, 2001.

\bibitem{gost19105}
ГОСТ 19.105-78: Общие требования к~программным документам. // Единая система программной документации.~--- М.: ИПК Издательство стандартов, 2001.

\bibitem{gost19106}
ГОСТ 19.106-78: Требования к~программным документам, выполненным печатным способом. // Единая система программной документации.~--- М.: ИПК Издательство стандартов, 2001.

\bibitem{gost19504}
ГОСТ 19.504-79: Руководство программиста. Требования к~содержанию и~оформлению. // Единая система программной документации.~--- М.: ИПК Издательство стандартов, 2001.

\bibitem{gost19604}
ГОСТ 19.604-78: Правила внесения изменений в~программные документы, выполненные печатным способом. // Единая система программной документации.~--- М.: ИПК Издательство стандартов, 2001.

\end{thebibliography}

\clearpage

% ============================================================
% ПРИЛОЖЕНИЕ 1 — ПРИМЕР КОМАНД ПРОГРАММИСТА
% ============================================================
\section*{ПРИЛОЖЕНИЕ 1}
\addcontentsline{toc}{section}{ПРИЛОЖЕНИЕ 1}

\begin{center}
\textbf{ПРИМЕР КОМАНД ПРОГРАММИСТА}
\end{center}

% Подсказка: приведите типовые команды программиста (запуск инфраструктуры,
% запуск сервисов, проверки, остановка окружения).
\begin{verbatim}
# Запуск инфраструктуры разработки
<команда запуска инфраструктуры>

# Запуск всех сервисов
<команда запуска сервисов>

# Выполнение проверок
<команда проверки>

# Остановка сервисов
<команда остановки>
\end{verbatim}

\clearpage

% ============================================================
% ЛИСТ РЕГИСТРАЦИИ ИЗМЕНЕНИЙ
% ============================================================
\input{../common/change_log}

\end{document}
